% Example content adapted from David Molik’s tenure document.
\PortfolioSection{Statement of Accomplishments}
\SectionGuidance{Summarize your most significant accomplishments during the applicable reporting period. Explain your role, outcomes, and significance across service and RSCAD, and point to supporting evidence.}
\ExampleNote

\subsection*{Directed Service}

\begin{itemize}
  \item Served as Department Head for the Center for Scholarly Publishing (renamed the Open Publishing Exchange), providing strategic and operational leadership across scholarly publishing, open educational resources, and research-support services.

  \item Presented the CSP/OPE landscape, scope, and strategic direction to the Associate Deans for Research Council (April 2025), strengthening integration of Libraries services within the University’s research enterprise.
    \begin{itemize}
      \item Authored a comprehensive departmental landscape document and executive summary (July 2025) that now serve as foundational planning and decision-making documents for CSP/OPE.
      \item Created and facilitated the ``Tea with CSP'' discussion series to support ongoing professional development and cross-campus engagement around scholarly communication, open publishing, and research infrastructure, and informing strategic priorities within the Open Publishing Exchange.
    \end{itemize}

  \item Oversaw early development of a Research Data Management and Curation program aligned with FAIR principles and federal open-access mandates (e.g., the Nelson Memo), including coordination with the Office of Research and Research IT.

  \item Oversaw migratory work from former New Prairie Press Digital Commons platform to Fulcrum/Janeway.

  \item Reorganized governance structures for the Open and Alternative Textbook Initiative (OATI), including drafting bylaws for a revised committee structure to improve transparency, shared ownership, and long-term sustainability (implementation scheduled for 2026).

  \item Participated in advanced leadership and assessment training through:
    \begin{itemize}
      \item Harvard’s Leadership Institute for Academic Librarians (LIAL).
      \item Kansas State University Fall New Department Head Institute.
      \item University of Kansas Assessment in Action Workshop.
    \end{itemize}

\end{itemize}

\subsection*{Non-Directed Service}

\textbf{Service to Kansas State University Libraries and the University}

\begin{itemize}

  \item Served as a member of the Environmental Health and Safety (EHS) Committee, contributing to institutional oversight related to research and workplace safety.

  \item Served on the Wildcat Scholar Committee, contributing scholarly communication and research infrastructure expertise to discussions related to research visibility, repository integration, and institutional research profiling.

  \item Taught at week-long Genomics Carpentries Workshop in August which was targeted towards incoming Life Sciences Students with Gwen Sibley and Kendra Spahr
\end{itemize}

\vspace{0.5em}

\textbf{Service to the Profession and National Communities}

\begin{itemize}
  \item Founded and served as Editor (and Co-Editor-in-Chief) of \textit{Cyberbiosecurity Quarterly}, a peer-reviewed scholarly journal hosted on New Prairie Press, establishing editorial governance, peer-review workflows, and a new scholarly venue at the intersection of genomics, cybersecurity, and biosecurity.

  \item Continued service with the i5k initiative, a global collaboration coordinating arthropod genome sequencing, annotation, and data standards, through participation in the Standards Working Group and by assuming leadership of the i5k Training Working Group, supporting community-wide best practices, workforce development, and the long-term usability and interoperability of genomic data resources.

  \item Served on the AgBioData Steering Committee and the AgBioData Scientific Literature Working Group, contributing to national coordination, governance, and standards development for agricultural genomics data. Began to serve as Co-Principal Investigator on the National Science Foundation Research Coordination Network project “Reimagining a Sustainable Data Network to Accelerate Agricultural Research and Discovery” (Award \#2126334), supporting annual community workshops, working group activities, and coordination among database, research, and scholarly communication stakeholders.

  \item Joined the advisory board of the Bioeconomy Information Sharing and Analysis Center (BIO-ISAC), which supports trusted information sharing and guidance on safeguarding biological data and research infrastructure, and participated in the BIO-ISAC annual conference, contributing to national conversations relevant to library responsibilities in data stewardship, scholarly communication, access, and institutional risk management.

  \item Served as Co-Principal Investigator with AgBioData, contributing to planning and coordination that brought the AgBioData annual meeting to Kansas State University Olathe. AgBioData is an NSF-funded Research Coordination Network of 45 Genomics, Genetics, and Breeding databases (GGBs) and Three publishers focused on sustaining and standardizing agricultural genomics data resources, work that closely aligns with library leadership in research data stewardship, interoperability, and long-term access. Bringing the AgBioData Annual Meeting to K-State Olathe elevates K-State’s visibility as a national hub for agricultural data infrastructure and positions the Libraries as a convening partner in research data stewardship.
\end{itemize}

\Redacted{Additional working updates}{Private consultation, proposal, and working-record details are omitted from this example.}

\subsection*{Research and Creative Activity}

\begin{itemize}
  \item Co-authored three peer-reviewed journal articles published in 2025:
    \begin{itemize}
      \item \textit{AgBioDatabase Finder: an online tool to help researchers find and submit agricultural genomic, genetic, and breeding data} (\textit{microPublication Biology}).
      \item \textit{BALROG-MON: a high-throughput pipeline for Bacterial AntimicrobiaL Resistance annOtation of Genomes--Metagenomic Oxford Nanopore} (\textit{microPublication Biology}).
      \item \textit{Guidelines for Gene and Genome Assembly Nomenclature} (\textit{Genetics}).
    \end{itemize}

  \item Contributed to these articles across the transition into my Kansas State University appointment, contributing to FAIR data practices, reproducible bioinformatics pipelines, and community standards for genomic data interoperability.

  \item Organized the session \textit{Cyberbiosecurity in Agricultural Genomics} and co-organized the session \textit{The AgBioData Consortium: Challenges and Recommendations for FAIR Genetic, Genomic, and Breeding Data} at the Plant and Animal Genome Conference (PAG33).

  \item Delivered the scholarly talk \textit{Some Thoughts on Genomics Data Integrity} at the \textit{Cyberbiosecurity in Agricultural Genomics} workshop at PAG33 (January 2026), contributing conceptual perspectives on data integrity, provenance, and trust in genomics.
\end{itemize}

\section{Accomplishments Evidence}
\SectionGuidance{Provide references or links supporting the accomplishments above. Make clear which claim each source supports.}
\printbibliography[keyword={achievements},heading=none]
